\section{Introduction}
\label{sec:introduction}

Relational event models (REMs) describe streams of time-stamped, directed
interactions --- emails between colleagues, messages in a classroom, citations
between documents, trips in a bike-sharing network --- as the realisations of a
continuous-time point process whose instantaneous rate depends on the history of
the network itself \citep{butts2008, bianchi2024}. A recent line of methodological work has extended this framework rapidly: timing,
closure and actor-heterogeneity effects \citep{juozaitiene2024}, models with
global covariates and a time-shift non-event sampling scheme
\citep{lembo2025}, martingale-residual goodness-of-fit diagnostics
\citep{boschi2025gof}, and relational \emph{hyper}-event models (RHEMs) with
time-varying non-linear effects \citep{boschi2025rhem, lerner2023}. Each
contribution is statistically self-contained and well-validated, but the
software that accompanies them is not: the methodology currently lives as
bespoke analysis scripts accompanying the individual papers, supplemented by
the Java tool \pkg{eventnet} for hyper-event feature extraction \citep{lerner2023}.
A practitioner who wants to simulate a network, sample non-events, compute
endogenous statistics, fit a model, and check its fit must today assemble these
stages by hand from heterogeneous sources.

The CRAN-native alternative does not close this gap either. The Tilburg
network group ships a mature, modular stack --- \pkg{remify}, \pkg{remstats},
\pkg{remstimate} for the analysis pipeline and \pkg{remulate} for simulation
\citep{lakdawala2025} --- but it is confined to \emph{linear} effects on the
log-rate and splits simulation from inference across separate packages with
distinct data contracts. The smooth, time-varying, and non-linear effects that
motivate much of this recent work, and the goodness-of-fit machinery
that accompanies them, have no home in that ecosystem. There is, in short, no
single tested R package that carries a relational (or hyper-relational) event
analysis from a simulated or observed event log all the way through to a fitted
model and its diagnostics.

\begin{table}[ht]
\centering\small
\begin{tabular}{lcccc}
\toprule
Capability & \pkg{amore} & \pkg{relevent} & \makecell{\pkg{remstats}\\stack$^{\dagger}$} & \makecell{\pkg{eventnet}\\(Java)} \\
\midrule
Linear endogenous effects        & \cmark & \cmark & \cmark & $\circ$ \\
Smooth / non-linear / time-varying & \cmark & \xmark & \xmark & \xmark \\
Random effects / frailty          & \cmark & $\circ$ & \xmark & \xmark \\
Flexible effects at scale (neural / SGD) & \cmark & \xmark & \xmark & \xmark \\
Event simulation                  & \cmark & \cmark & $\circ$ & \xmark \\
Case-control non-event sampling   & \cmark & \xmark & $\circ$ & \cmark \\
Residual-based goodness of fit    & \cmark & \xmark & \xmark & \xmark \\
Hyper-event (set-valued) support  & \cmark & \xmark & \xmark & \cmark \\
Single package, one data contract & \cmark & $\circ$ & \xmark & \xmark \\
\bottomrule
\end{tabular}
\caption{Primary documented scope of relational-event software, relative to
\pkg{amore}. \cmark: supported in the package; $\circ$: supported in
restricted form, in a companion package, or as an external tool; \xmark:
absent. The comparison reflects the packages' documented capabilities as of
the cited versions and is necessarily coarse. $^{\dagger}$The Tilburg stack
(\pkg{remify}, \pkg{remstats}, \pkg{remstimate}, \pkg{remulate};
\citealp{lakdawala2025}) is modular: simulation (\pkg{remulate}) and inference
are separate packages with distinct data contracts, and the estimation
pipeline is restricted to linear effects on the log-rate.
\pkg{eventnet} \citep{lerner2023} is a Java feature-extraction tool, including
hyper-event statistics and nested case-control sampling, paired with external
model fitting. \pkg{relevent} \citep{butts2008} is the original \proglang{R}
REM implementation, covering linear models and simulation.}
\label{tab:landscape}
\end{table}

Table~\ref{tab:landscape} situates \pkg{amore} against the main alternatives.
No existing tool offers smooth, neural, and goodness-of-fit machinery within a
single simulation-to-diagnostics workflow; that combination, rather than any
individual capability, is what \pkg{amore} contributes.

\subsection{What \pkg{amore} provides}
\label{sec:what-amore-provides}

\pkg{amore} closes that loop: a single \proglang{R} package in which this
strand of relational event methodology becomes one coherent workflow
spanning five stages:

\begin{enumerate}
  \item \emph{simulation} of event streams by an exact Gillespie algorithm or an
  approximate tau-leap scheme sharing one endogenous-state machine
  (\code{simulate\_relational\_events});
  \item \emph{case-control non-event sampling}, including the \citet{lembo2025}
  time-shift scheme for global covariates, with sampling scopes for repeatable
  and non-repeatable processes (\code{sample\_non\_events});
  \item a history-aware \emph{endogenous-statistics engine} covering the
  full catalogue of timing and closure variants
  (\code{compute\_endogenous\_features});
  \item a single \emph{model-fitting} front-end, \code{rem()}; and
  \item the \citet{boschi2025gof} \emph{martingale-residual goodness-of-fit}
  family for the linear case-control fit (\code{gof\_univariate},
  \code{gof\_multivariate}, \code{gof\_global}, \code{gof\_auxiliary}).
\end{enumerate}

The five stages share one data contract --- the \code{amore\_event\_log} ---
so that the output of the simulator feeds the sampler, whose output feeds the
feature engine, whose output feeds \code{rem()} directly. Correctness across
the seam between simulation and post-hoc analysis is enforced by a
cross-implementation parity oracle (the simulator's
on-the-fly statistics are compared against the independently coded post-hoc
engine to a tolerance of $4.3\times10^{-14}$ under aligned histories), backed
by an extensive automated test suite.

\subsection{One likelihood, a spectrum of predictors}
\label{sec:one-likelihood}

The expository spine of the package is a single estimand viewed through three
fitting backends, selected by \code{rem(method = c("gam", "clogit", "nn"))}.
Two of these maximise \emph{the same} objective: the risk-set softmax
conditional-logistic partial likelihood that arises when the continuous-time
REM intensity is reduced, via case-control sampling, to a conditional
multinomial choice among the realised event and its sampled non-events
\citep{lerner2020}. The \code{clogit} backend fits this partial likelihood with
a linear predictor through \pkg{survival}; the \code{nn} backend maximises the
identical per-stratum objective with a multilayer-perceptron predictor and
hand-derived analytic gradients. The third backend, \code{gam}, fits the
\emph{algebraically related} degenerate case-1-control binomial logistic
regression of \citet{boschi2025gof, juozaitiene2024} through \pkg{mgcv}, which
unlocks smooth, time-varying, non-linear, and random-effect terms. We are
deliberately precise about this distinction: the \code{gam} backend is not a
risk-set softmax fit and coincides with a two-element conditional logit only
under a single matched control, a condition \code{rem()} does not enforce.
What \code{rem()} offers is therefore not one likelihood maximised three ways,
but a single front-end and a shared \code{rem} S3 contract spanning a
spectrum of predictors --- linear, smooth, and neural --- over closely related
case-control objectives. To our knowledge no prior work presents this
linear-to-smooth-to-neural progression as one workflow behind one interface.

\subsection{What this paper does not claim}
\label{sec:scope}

Because \pkg{amore} sits squarely inside an active methodological literature,
we state its boundaries plainly. \emph{The contribution is software and
synthesis, not new statistics.} Every estimator, sampling scheme, and
diagnostic implemented here is published or derives directly from the
published literature; \pkg{amore}'s value is the integrated, documented, tested, and
reproducible artifact, together with the unifying frame of
Section~\ref{sec:one-likelihood}.

In particular, the \code{nn} backend is \emph{not} a claim to the first neural
REM. Flexible-effect relational event modelling at scale already exists:
STREAM \citep{filippimazzola2024} fits additive B-spline effects by minibatch
stochastic gradient descent on the same nested case-control objective, scaling
to $10^{8}$ patent citations, and the neural temporal-point-process literature
--- recurrent marked temporal point processes \citep{du2016}, the neural
Hawkes process \citep{mei2017}, dynamic-graph representation learning
\citep{trivedi2019}, and the review by \citet{shchur2021} --- long predates
this package. What
\pkg{amore}'s \code{nn} backend contributes is accessibility and engineering: a
GPU-free, zero-extra-dependency, pure-R neural conditional-logit backend,
reachable by \code{install.packages()}, whose analytic gradients are
verified against numerical differentiation (agreement to
$1.4\times10^{-10}$), and which --- by scoring the full candidate vector
jointly rather than additively --- can learn cross-feature interactions and
recover them post-hoc through partial-dependence plots. Its reported
concordance figures are \emph{in-sample}; it provides no uncertainty
quantification (\code{coef()} returns \code{NULL}, there are no standard errors
or variance-covariance estimates), it trains by gradient descent in pure R
(full-batch by default, with optional mini-batching), and it does not accept
factor covariates. We make no scalability or inferential claim
for it. Finally, the goodness-of-fit family is restricted to the linear
case-control fit (\code{rem()} rejects non-linear specifications before any GOF
computation), and the RHEM hyper-event statistics, while \code{rem}-ready, are
not yet demonstrated end-to-end on a fitted hyper-event model.

\subsection{Contributions}
\label{sec:contributions}

We organise the contribution in three tiers, from the core deliverable to
supporting material.

\paragraph{Primary: software and synthesis.}
\pkg{amore} is, to our knowledge, the first \proglang{R} package that closes
the full REM/RHEM loop --- simulate, sample non-events, compute features,
fit, and check --- for these recent formulations
\citep{juozaitiene2024, lembo2025, boschi2025gof, boschi2025rhem}, which
have so far been available only as scattered analysis scripts and the Java
\pkg{eventnet} tool, while existing CRAN packages cover linear effects and
split simulation from inference. The accompanying expository frame --- one
front-end and one S3 contract over a linear-to-smooth-to-neural spectrum of
predictors built on closely related case-control objectives
(Section~\ref{sec:one-likelihood}) --- is, as far as we are aware, not
articulated as a single workflow elsewhere.

\paragraph{Software contributions.}
The package ships with strong validation evidence, including a genuine
cross-implementation oracle: simulator-versus-post-hoc feature parity to
$4.3\times10^{-14}$ under aligned histories, planted-effect recovery under both
simulation schemes, an analytic-versus-numerical gradient check for the neural
backend ($1.4\times10^{-10}$), and the goodness-of-fit identities (including an
analytic single-component reduction verified to $10^{-10}$). The
pure-R neural conditional-logit backend, described above, is a self-contained
contribution: analytic gradients verified against numerical differentiation,
trained on the exact case-control partial likelihood, with no external
numerical dependency.

\paragraph{Supporting.}
\pkg{amore} bundles a curated battery of 9 documented data objects spanning
roughly two orders of magnitude in event count (439 to 82{,}927 events), with
source DOIs and rebuild scripts, giving the framework a graded, reproducible
stress test across diverse directed one-mode interaction networks --- email,
online messaging, phone calls, classroom interaction, and friendship.

\subsection{Roadmap}
\label{sec:roadmap}

Section~\ref{sec:background} reviews the REM likelihood, its reduction to a
case-control softmax over sampled non-events, the endogenous-statistics
catalogue, and the time-shift device for global covariates, fixing the
estimands the backends share. Section~\ref{sec:framework} treats the three
backends and the precise sense in which they share an objective.
Section~\ref{sec:software} covers the software design --- the case-control
data pipeline (event-log canonicalisation, non-event sampling, feature
computation, and the wide case-control representation), the \code{rem()} S3
contract, the goodness-of-fit family, the bundled datasets, and a worked
end-to-end example. Section~\ref{sec:simulation} describes the Gillespie and
tau-leap simulation engine. Section~\ref{sec:validation} reports the
validation experiments E1--E6, and Section~\ref{sec:illustration} works
through illustrations on the bundled real datasets.
Section~\ref{sec:discussion} discusses limitations and outlines
the road ahead --- an end-to-end RHEM demonstration, a minibatched neural
backend with uncertainty quantification, and goodness-of-fit beyond the linear
fit.
